\chapter{Discussion}\label{ch:discussion}

\section{What the baseline establishes}
The baseline demonstrates that a physics-informed field representation can
agree closely with an independently discretized curved-channel solution
for the specified steady, low-Reynolds-number problem. The velocity and
pressure errors are below 2\% and 1\%, respectively, while the drag-network
error is below 1\%. The analytical plane-flow test provides an additional
check that is independent of the curved-channel finite-element field.
Reference-mesh sensitivity and section flux add evidence beyond a single
parity plot.

The trajectory comparison evaluates an additional quantity beyond field
agreement. Six of seven endpoint classes agree, while three of seven
trajectories meet the complete position, common-time and endpoint criteria.
A particle integrates the local field over time; force approximation,
interpolation and event handling also affect its result. Their individual
contributions require a controlled comparison. These findings motivate
direct transport assessment alongside the carrier-field metrics.

The force network also raises a methodological question. Since the
retained Stokes law is known analytically, its neural approximation is
most useful as a controlled component of a broader learning framework.
The present error does not show that the learned force is superior to the
direct analytical formula. A future assessment should compare the two
within identical carrier fields and identify whether the approximation
adds cost, flexibility, or error.

\section{What conservative transfer establishes}
The shared Gaussian operator gives a stronger result of a different kind.
The force and work identities follow from the normalization and the common
quadrature. Their validity does not depend on a low training loss. This
property removes one avoidable source of inconsistent interphase exchange
and remains available when the carrier geometry changes, provided that the
same fluid-domain quadrature is used consistently.

The identities do not guarantee the accuracy of the coupled solution.
The carrier momentum equation can retain a residual even when the sum of
deposited forces is exact. In the evaluated coupled channel, the remaining
momentum residual exceeds the study's bound in every tested parameter
combination. The correct inference is that the operator design is verified
while the quantitative coupled-channel prediction remains unresolved.

The successful reduced tests address important parts of the calculation:
response-time dynamics, reciprocal momentum exchange, kernel integrals,
and unsteady incompressible flow. Each isolates a necessary capability.
Their combination is a verification chain, not a substitute for a matched
independent two-way channel reference. This distinction is consistent with
the separation between numerical verification and physical validation
\citep{roache1998}.

\section{Optimization, residuals, and independent error}
The completed calculations illustrate why optimization metrics and
independent solution errors should be reported together. Residual
normalization changes their numerical magnitude, and a weighted objective
can favor some components over others. Reported raw residuals should
therefore retain their equation convention. In particular, comparing the
baseline momentum RMS directly with the Gaussian momentum RMS would
ignore different scaling, geometry, and transient content.

Training difficulties in physics-informed learning are well documented
\citep{wang2021,wang2022,krishnapriyan2021}. They motivate focused improvements such as % Citation-plan C31
component scaling and more informative collocation. Residual-based
adaptive sampling \citep{wu2023} and adaptive loss weighting
\citep{mcclenny2023} are options for a future controlled comparison.
A further option is better temporal % Citation-plan C32
resolution of the response scale, including causality-aware training
for the transient coupled problem \citep{wang2024}. These approaches were
not part of the recorded runs. Their value must be assessed % Citation-plan C33
using the same independent checks and the original acceptance criteria.

\section{Regularization and model validity}
The Gaussian width controls how the represented force is distributed.
It affects both the sampled fluid velocity and the carrier response.
A smaller width is not automatically more physically accurate, especially
if the transfer quadrature underresolves it. A larger width can smooth
localized forcing while changing the regularized interaction. Width
sensitivity must therefore be distinguished from convergence at fixed
model parameters.

The undisturbed-velocity issue is equally important. The current parcels
sample the disturbed fluid directly \citep{horwitz2016,ireland2017}. % Citation-plan C34
 A reduced Stokes-response test cannot
resolve the resulting self-interaction question in a spatially coupled
flow. Comparison with an independent solver of the same regularized
equations would test numerical consistency; assessment against a suitable
physical reference would additionally test the closure.

Confinement further separates numerical and physical accuracy. A model
can converge accurately to equations that omit important wall effects.
For that reason, the proposed constriction study records particle-to-gap
ratio and local slip Reynolds number in addition to ordinary solver errors.
If those quantities indicate a breakdown of the retained approximation,
more training is not the appropriate physical correction.

\section{Transferability across the three stages}
\Cref{tab:claim_scope} summarizes the actual scope of the findings. The
baseline establishes a carrier-field comparison for one geometry and
Reynolds number, with two seeds. The Gaussian tests establish specific
reduced dynamics and transfer properties. The constriction stage provides
a defined study design. None of these establishes a universal neural
operator over channel shape, loading, and Reynolds number.

\begin{table}[htbp]\centering\small
\caption{Supported conclusions and their boundaries.}\label{tab:claim_scope}
\begin{tabularx}{\textwidth}{@{}p{.24\textwidth}YY@{}}\toprule
Stage & Supported statement & Remaining boundary\\\midrule
Baseline carrier & Analytical and independent numerical field agreement & One declared geometry/Re case; finite residuals\\
Baseline particles & Completed integrations and drag approximation & Six of seven matching endpoint classes; three complete path comparisons within criteria\\
Gaussian transfer & Discrete force/work consistency and component verification & Disturbed-velocity closure and coupled field error\\
Straight coupled channel & Executed diagnostic predictions & Momentum RMS $0.0407$--$0.3854$ against a $0.03$ criterion\\
Constriction & Defined geometry and evaluation protocol & No completed constricted-channel solution\\
\bottomrule\end{tabularx}
\end{table}

The transferable parts are the nondimensional bookkeeping, reference
matching, shared transfer, and evaluation logic. New physical and geometric
claims require their own executed evidence. In particular, a parameter
slot in a network does not establish learned parametric generalization,
and a geometric illustration does not establish a flow solution.

\section{Implications for the next computational study}
The most informative next work is narrower than a large parameter sweep.
First, the baseline trajectory limitation should be isolated using the
same carrier field with exact and learned drag. Second, the Gaussian
channel should be compared with an independent regularized solver at fixed
width and well-resolved quadrature. Third, the unladen constricted problem
should be solved before introducing moving populations.

This ordering separates carrier approximation, irregular-domain transfer,
wall geometry, particle closure, and optimization. Establishing the controls
in sequence allows each additional mechanism to be interpreted against a
known numerical foundation.

\section{Relationship to earlier methods}
The present contribution is a continuous-time neural representation with
shared, normalized force transfer and explicitly separated field and
trajectory assessment. The transfer structure is related to
Eulerian--Lagrangian methods \citep{peskin2002,capecelatro2013}, while the
fluid representation belongs to the velocity--pressure PINN framework
\citep{jin2021}. Undisturbed-velocity recovery is a distinct modeling
development addressed by Horwitz and Mani \citep{horwitz2016} and Ireland
and Desjardins \citep{ireland2017}. These
relationships identify the design choices adopted here; the reported
comparisons establish performance only for the stated tests and do not
demonstrate general superiority over those earlier methods.
